% from ega0-18-fr.tex
\oldpage[0\textsubscript{IV-1}]{57}
$\theta_2=u\circ\theta_1$~; cela résulte aussitôt de ce qu'on a vu dans
(18.2.5) et de la définition de la bijection canonique de $F$ sur le produit
fibré $E'\times_{B'}B$ (18.1.7).
\end{env}

\begin{env}[18.2.7]
\label{0.18.2.7-fr}
Considérons maintenant le carré de gauche de (18.2.4.1), et rappelons d'abord
la notion de \emph{somme amalgamée} dans la catégorie des $A$-bimodules~:
étant donnés trois $A$-bimodules $X$, $Y$, $Z$ et deux $A$-homomorphismes
$f:X\to Y$, $g:X\to Z$, la somme amalgamée $Y\oplus_XZ$ est \emph{limite
inductive} du système inductif formé par $X$, $Y$, $Z$ et les
$A$-homomorphismes $f$, $g$ dans la catégorie des $A$-bimodules
($0_{\mathrm{III}}$, 8.1.11). On définit ce $A$-bimodule comme quotient du
produit $Y\times Z$ par le sous-$A$-bimodule $M$ image de $X$ par
l'homomorphisme $x\mapsto(f(x),-g(x))$. Sa propriété caractéristique est que,
pour tout couple d'homomorphismes de $A$-bimodules $u:Y\to T$, $v:Z\to T$
tels que $u\circ f=v\circ g$, il existe un homomorphisme et un seul
$w:Y\oplus_XZ\to T$ tel que $u=w\circ j_1$ et $v=w\circ j_2$, où
$j_1:Y\to Y\oplus_XZ$ et $j_2:Z\to Y\oplus_XZ$ sont les applications
canoniques.
\end{env}

\begin{env}[18.2.8]
\label{0.18.2.8-fr}
Considérons maintenant une $A$-extension $E$ de $B$ par un $B$-bimodule
$L$~:
\[
  0\longrightarrow L\xrightarrow{j}E\xrightarrow{f}B\longrightarrow0
\]
et soit d'autre part $w:L\to L'$ un homomorphisme de $B$-bimodules. Soit $H$
le $A$-bimodule somme amalgamée $E\oplus_LL'$~; montrons comment on peut munir
ce $A$-bimodule d'une structure de $A$-anneau et définir une $A$-extension
\[
  0\longrightarrow L'\longrightarrow H\longrightarrow B\longrightarrow0.
\]

Notons pour cela que $L'$ est muni d'une structure de $E$-\emph{bimodule} au
moyen de l'homomorphisme d'augmentation $E\to B$~; on peut donc former la
$A$-extension triviale type $G=D_E(L')$ (18.2.3). Considérons alors
l'application $\theta:z\mapsto(j(z),-w(z))$ de $L$ dans $G$~; c'est un
homomorphisme de $G$-bimodules ($L$ étant considéré comme $G$-bimodule au
moyen de l'homomorphisme canonique $\rho:G\to E$). En effet, pour
$(x,z')\in G$ et $z\in L$, on a
$(j(z),-w(z))(x,z')=(j(z)x,j(z)z'-w(z)f(x))$~; or
$j(z)z'=f(j(z))z'=0$ par définition de la structure de $E$-bimodule sur
$L'$, et $w(z)f(x)=w(zf(x))$ par définition de la structure de $B$-bimodule
sur $L'$~; on vérifie de même que $\theta$ est un homomorphisme de
$G$-module à gauche. On peut alors appliquer au diagramme commutatif
\[
  \xymatrix@C=2pc@R=1.4pc{
    & & 0\ar[d] & & \\
    & & L\ar[ld]^\theta\ar[d]^j & & \\
    0\ar[r] & L'\ar[r] & G\ar[r]_p & E\ar[r]\ar[d]^f & 0\\
    & & & B\ar[d] & \\
    & & & 0 &
  }
\]
le résultat de (18.1.7). Comme $H=G/\theta(L)$ par définition, notre
assertion est une conséquence immédiate de (18.1.7).

